\scanpage{31}
% L1–4
Soit $f:\deleted{\unsure{Y}}X\longrightarrow\deleted{\unsure{X}}Y$ un $S$-morphisme. On appelle \uline{transformation infinitésimale de type} $S'$ \uline{de} $f$ un couple $u,v$ de $S'$-automorphismes de $X_{S'},Y_{S'}$ \inserted{induisant l'identité sur $X,Y$} rendant commutatif le diagramme
% L5–6
\diagram{\begin{tikzcd}[ampersand replacement=\&]
X_{S'}\arrow[r,"u"]\arrow[d,"f_{S'}"']\&X_{S'}\arrow[d,"f_{S'}"]\\
Y_{S'}\arrow[r,"v"]\&Y_{S'}
\end{tikzcd}}
% L7–9
Les transformations infinitésimales \inserted{$w=(u,v)$} de type $S'$ de $f$ forment un \uline{groupe}, contenu dans le produit des groupes des autom. inf. de type $S'$ de $X$ et $Y$.

% L10–13
Soit $w=(u,v)$ une transformation infinitésimale de type $S'$ de $f$, et soit $g$ une section de $X/Y$. Considérons alors
\[
\widetilde w(g)=u\,g_{S'}\,v^{-1}
\]
% L14–20
c'est une section de $X_{S'}$ sur $Y_{S'}$ ``proche'' de $g$, donc correspond à \unsure{une} \deleted{\unsure{S-flèche|flèche|système} $\deleted{\unsure{X}}Y\times_S S'{\to}X$ \unsure{la} \unsure{flèche} tel que son composé avec $X\unsure{-|\to}Y$ soit $\operatorname{pr}_1$, i.e. un $S$-hom $Y\to X^{\{S'\}}$ tel que son composé avec $X^{\{S'\}}\unsure{-|\to}Y^{\{S'\}}$ soit la section \unsure{évidente} $Y\to Y^{\{S'\}}$ \unsure{section de}} élément de $\Gamma(X_Y^{\{Y_{S'}\}}/Y)$ ``proche'' de $g$.

% L21
Propriétés fonctorielles triviales \ldots\ldots
% archive sheet number: 30

